\begin{figure*}[htbp]
\includegraphics[width=\textwidth]{rich_mode_frame}
\caption{Rich mode: two-column layout with IR source spotlight alongside the stack.}
\label{fig:rich_mode}
\end{figure*}

LLVManim provides three stack-focused animation modes and one graph-based mode, each reproducible with a single CLI command. The stack-focused modes (basic, rich, rich-SSA) share a common layout and animation style, with increasing levels of IR detail. The graph-based mode renders an animated traversal of the control-flow graph using Graphviz layout and a runtime trace for the traversal order. Each mode is designed to target different presentation contexts and audience familiarity levels, allowing users to choose the appropriate level of detail for their needs. The following sections describe each mode in detail, along with example CLI commands for generating the animations. All animations were generated from the same \texttt{double.ll} test program, which contains a simple loop with arithmetic operations, memory accesses, and control flow. The full source code for \texttt{double.c} and \texttt{double.ll} is provided in appendix~\ref{app:code_examples}.

\vspace{2ex}

\textbf{Basic mode} (\texttt{-{}-ir-mode basic}) renders a minimal stack-only view. Each IR instruction triggers a per-cell badge flash on the corresponding stack slot, providing a clean visual pulse of execution without additional clutter. This mode is best suited for audiences who need only a high-level sense of execution flow, such as introductory lectures or quick demonstrations, since it strips away all source-level detail and focuses purely on the call stack. Listing~\ref{cli:basic_mode} shows the CLI command used to generate the animation shown in Fig.~\ref{fig:basic_mode}.
\vspace{-4ex}
\begin{minted}[
  autogobble,
  frame=single,
  bgcolor=listingbg,
  breaklines=true,
  fontsize=\footnotesize,
]{bash}
  uv run llvmanim double.ll --animate --ir-mode basic --outdir out/
\end{minted}
\vspace{-5ex}
\captionof{listing}{CLI command to generate the basic mode animation shown in Fig.~\ref{fig:basic_mode}.}
\label{cli:basic_mode}

\vspace{2ex}

% MOVED FIGURE TO RISK_EVALUATION.TEX SO IT WOULD RENDER ON THE SAME PAGE --- IGNORE ---
% \begin{figure*}[htbp]
% \includegraphics[width=\textwidth]{basic_mode_frame}
% \caption{Basic mode: stack-only view with per-cell badge flash during a \texttt{call} to \texttt{@init}.}
% \label{fig:basic_mode}
% \end{figure*}

\textbf{Rich mode} (\texttt{-{}-ir-mode rich}) extends the basic view with a two-column layout. An IR source panel on the left displays the current function's instructions while a spotlight cursor tracks the executing line, giving the audience both the ``what'' (the stack) and the ``where'' (the source). This mode is particularly valuable when the presenter wants to walk through specific IR constructs---for example, showing how a series of \texttt{alloca} and \texttt{store} instructions initialize local variables before a loop body executes. The spotlight cursor provides a natural focal point that guides the viewer's attention through the code. Listing~\ref{cli:rich_mode} shows the CLI command used to generate the animation shown in Fig.~\ref{fig:rich_mode}.
\vspace{-2ex}
\begin{minted}[
  autogobble,
  frame=single,
  bgcolor=listingbg,
  breaklines=true,
  fontsize=\footnotesize,
]{bash}
  uv run llvmanim double.ll --animate --ir-mode rich --outdir out/
\end{minted}
\vspace{-5ex}
\caption{CLI command to generate a rich mode animation shown in Fig.~\ref{fig:rich_mode}.}
\label{cli:rich_mode}
\vspace{2ex}

\textbf{Rich-SSA mode} (\texttt{-{}-ir-mode rich-ssa}) takes the rich layout a step further with a dedicated third column for SSA values (IR Source $\mid$ SSA Values $\mid$ Stack). The center panel displays the symbolic result of each binary operation, comparison, or load as it executes, making it possible to follow data dependencies without mentally resolving \texttt{\%}-names. This is the most detailed mode and targets audiences already familiar with SSA form---for instance, a compiler course where students need to see how \texttt{\%4 = icmp slt i32 \%2, \%3} evaluates during a loop iteration. Listing~\ref{cli:rich_ssa_mode} shows the CLI command used to generate the animation shown in Fig.~\ref{fig:rich_ssa_mode}.
\vspace{-2ex}
\begin{minted}[
  autogobble,
  frame=single,
  bgcolor=listingbg,
  breaklines=true,
  fontsize=\footnotesize,
]{bash}
uv run llvmanim double.ll --animate --ir-mode rich-ssa --outdir out/
\end{minted}
\vspace{-5ex}
\caption{CLI command to generate a rich-SSA mode animation shown in Fig.~\ref{fig:rich_ssa_mode}.}
\label{cli:rich_ssa_mode}
\vspace{2ex}

\begin{figure*}[htbp]
\includegraphics[width=\textwidth]{rich_ssa_mode_frame}
\caption{Rich-SSA mode: three-column layout showing computed SSA values for \texttt{load} and \texttt{icmp} instructions.}
\label{fig:rich_ssa_mode}
\end{figure*}

\needspace{6\baselineskip}
\textbf{CFG traversal} (\texttt{-{}-cfg-animate}) renders an animated walk through the control-flow graph, using Graphviz DOT layout for node positions and a runtime trace for the traversal order:
\vspace{-2ex}
\begin{minted}[
  autogobble,
  frame=single,
  bgcolor=listingbg,
  breaklines=true,
  fontsize=\footnotesize,
]{bash}
  uv run llvmanim double.ll --cfg-animate --dot-cfg main.dot --import-trace trace.json --outdir out/
\end{minted}
\vspace{-5ex}
\caption{CLI command to generate a CFG traversal animation.}
\label{cli:cfg_traversal}
\vspace{2ex}
